% !PRINTED 484
\weberpage{484}
%<<CONT
also:
\[ f_2\left(\frac{\omega}{n}\right) = e^{-\frac{r\pi i}{24}}x, \]
\[ f_2(n\omega) = \frac{e^{-\frac{r\pi i}{24}}\sqrt2}{x}, \]
wenn, je nachdem $r$ gerade oder ungerade ist,
\weq{33}{x = f_1(\sqrt{-m}) \quad\text{oder}\quad = f(\sqrt{-m})}
% [sic] printed omits $x$ before the second equality
gesetzt wird. Nehmen wir also in der zu $n$ gehörigen Modulargleichung
 II., § 73:
\[ u_1 = f_2(\omega), \]
\[ v_1 = \left(\frac2n\right) f_2(n\omega) = \left(\frac2n\right) \frac{e^{-\frac{r\pi i}{24}}\sqrt2}{x} \]
oder
\[ v_1 = f_2\left(\frac{\omega}{n}\right) = e^{-\frac{r\pi i}{24}}x, \]
so ergeben sich zwei Gleichungen, aus denen man durch Elimination
 von $u_1$ eine Gleichung für $x$ herleitet.
Um für $n = 5$ diese Rechnung durchzuführen, setzen wir:
\weq{34}{\begin{aligned}
x f_2(\omega) &= \sqrt2\,e^{-\frac{5r\pi i}{24}}\xi, \\
\frac{x}{f_2(\omega)} &= e^{-\frac{5r\pi i}{24}}\eta,
\end{aligned}}
woraus man, entsprechend den beiden Annahmen für $v_1$, für ein
ungerades $r$ die Gleichungen erhält:
\weq{35}{\begin{aligned}
\xi^3+\frac1{\xi^3}+2\left(\eta^2-\frac1{\eta^2}\right)&=0, \\
\eta^3+\frac1{\eta^3}+2\left(\xi^2-\frac1{\xi^2}\right)&=0.
\end{aligned}}
Diese Gleichungen gelten für $r = 1, 3, 7$, die zu den Werten
$m = 49, 41, 1$ gehören.
Subtrahiert man die beiden Gleichungen (35), so kann man
den Faktor $\xi - \eta$ abwerfen, der nur für $m = 1$ verschwinden
kann:
\[ [(\xi+\eta)^2 - \xi\eta]\left(1-\frac1{\xi^3\eta^3}\right) - 2(\xi+\eta)\left(1+\frac1{\xi^2\eta^2}\right) = 0, \]
%CONT>>
